%%%%%%%%%%%%%%%%%%%%%%%%%%%%%%%%%%%%%%%%%%%%%%%%%%%%%%%%%%%%%%%%%%%%%%%%%%%%%%%%
%% CHAPTER 2 SECTION: DETAILED LITERATURE REVIEW: AGENTIC AI DISEASE FINDER
%%%%%%%%%%%%%%%%%%%%%%%%%%%%%%%%%%%%%%%%%%%%%%%%%%%%%%%%%%%%%%%%%%%%%%%%%%%%%%%%

\section{Detailed Literature Review: Agentic AI Disease Finder}
\label{app:ch2_agentic}
\label{sec:topic5_agentic}

This section provides the detailed literature review for multi-modal medical AI and the Agentic Disease Finder (NeuroMCP-Agent) supporting the thematic synthesis in Chapter~2. The review covers the evolution from single-modality deep learning through multi-modal fusion strategies to foundation models and agentic architectures with intelligent model selection. All citations are referenced in the consolidated bibliography.

The literature on multi-modal medical AI has evolved from single-modality deep learning (2018--2020), through multi-modal fusion strategies (2020--2022), to foundation models (2022--2024) and agentic architectures with intelligent model selection (2024--2025).

\subsection{Agentic and Multi-Modal Systems}

Four multi-agent architectures define the current landscape, but each exhibits a fundamental limitation for biomedical deployment. Zhou et al.~\cite{zhou2025mam} rely entirely on LLM-based agents without domain-specific trained models, risking hallucinated medical reasoning. Jin et al.~\cite{jin2025medagent} add evidence-based case reasoning (MedAgent-Pro) but cannot route EEG or MRI data to specialised classifiers. Wang et al.~\cite{wang2025llmagent} combine LLM reasoning with medical imaging but lack real-time signal processing capability. Chen et al.~\cite{chen2025adacomed} employ an always-fuse strategy that cannot selectively engage the appropriate expert, imposing unnecessary computational overhead. The absence of EEG-specific routing in all four systems leaves a gap for ASD-focused diagnostic orchestration.

\subsection{Motor Imagery BCI}

Wang et al.~\cite{wang2024misurvey} comprehensively reviewed CNNs, RNNs, transformers, and hybrid models on BCI Competition IV-2a, where best models reach 80--85\% accuracy. Gu et al.~\cite{gu2025cltnet} achieved $\sim$83\% accuracy with a CNN + LSTM + transformer hybrid.

\subsection{PD Detection}

Jibon et al.~\cite{jibon2024hybrid} improved PD detection accuracy using a hybrid autoencoder-CNN model. Rabie et al.~\cite{rabie2025pd} surveyed DL methods from speech, gait, and EEG modalities for PD detection and monitoring.

\subsection{MRI Tumor Classification}

Zahoor et al.~\cite{zahoor2024resbrnet} achieved $>$95\% accuracy with a deep residual region-based CNN (Res-BRNet) combining global and local features. Srinivasan et al.~\cite{srinivasan2024cnn} compared three fully automatic CNN models for multi-class tumor classification.

\subsection{Agentic Innovation}

The proposed NeuroMCP-Agent uniquely combines concrete specialized expert models (BCI2A, PD EEG, MRI Tumor) with intelligent agentic routing based on input characteristics. The system achieves: PD detection at 100\% accuracy (5-fold CV), epilepsy detection at 99.02\%, ASD detection at 97.67\%, and an overall mean accuracy of 96.14\% across seven diseases. The system deploys 15 base classifiers with an MLP meta-learner, incorporates 46 Responsible AI modules, and achieves regulatory compliance scores of EU AI Act 95\%, FDA SaMD 94\%, and HIPAA 94.2\%.

\subsection{Comparative Analysis}

Table~\ref{tab:agentic_comparison} compares existing multi-modal medical AI systems against the proposed Agentic Disease Finder.
\needspace{8\baselineskip}
{\tablestripes\tablefontsize
\begin{xltabular}{\textwidth}{@{} p{3cm} p{3cm} p{2.5cm} L L @{}}
\caption[Comparative Analysis of Multi-Modal Medical AI vs.\]{Comparative Analysis of Multi-Modal Medical AI vs.\ Agentic Disease Finder}\label{tab:agentic_comparison} \\
\toprule
\thead{Study} & \thead{Method} & \thead{Result} & \thead{Gap} & \thead{Agentic Advantage} \\
\midrule
\endfirsthead
\multicolumn{5}{@{}l}{\textit{(\,Tab.~\ref{tab:agentic_comparison} continued from previous page\,)}} \\
\toprule
\thead{Study} & \thead{Method} & \thead{Result} & \thead{Gap} & \thead{Agentic Advantage} \\
\midrule
\endhead
\bottomrule
\endlastfoot
Zhou et al.~\cite{zhou2025mam} & MAM modular multi-agent & Outperforms baselines & LLM-centric; no EEG/MRI experts & Concrete EEG/MRI expert models + routing \\
\addlinespace
Jin et al.~\cite{jin2025medagent} & MedAgent-Pro evidence-based & Superior reliability & Imaging/text only; no EEG & Three-domain coverage with EEG + MRI \\
\addlinespace
Chen et al.~\cite{chen2025adacomed} & AdaCoMed fusion & Strong collaborative fusion & Always fuses; no expert selection & Input-aware model selection via agent \\
\addlinespace
Jibon et al.~\cite{jibon2024hybrid} & Hybrid AE + CNN for PD & Improved PD detection & Single disease, single modality & Multi-disease with agentic routing \\
\addlinespace
Zahoor et al.~\cite{zahoor2024resbrnet} & Res-BRNet for brain tumor MRI & $>$95\% accuracy & Tumor imaging only; no EEG & Unified architecture spanning EEG + MRI \\
\end{xltabular}}
\vspace{0.5em}\noindent\textit{Note.} MAM = modular multi-agent; LLM = large language model; AE = autoencoder; CNN = convolutional neural network; MRI = magnetic resonance imaging.
